\chapter{Resultados}

Este capítulo apresenta a análise e a discussão dos resultados obtidos a partir do desenvolvimento e da execução da solução computacional proposta, avaliando o comportamento da arquitetura sob carga e a eficácia das regras de negócio metrológicas implementadas.

\section{Desempenho e Escalabilidade da Arquitetura}

\textcolor{red}{(Apresentar aqui os resultados dos testes de estresse e desempenho dos microsserviços conteinerizados sob arquitetura orientada a eventos. Incluir gráficos e dados estatísticos que comprovem a capacidade de processamento assíncrono em lotes volumosos de certificados, taxas de transferência de mensagens no barramento, tempos médios de resposta e a eficiência da replicação de contêineres sem degradação da base transacional.)}

\section{Eficácia do Batimento Automatizado e Alertas}

\textcolor{red}{(Demonstrar aqui as métricas e os resultados obtidos no algoritmo de batimento entre a base interna e os relatórios/certificados enviados pelo Mock da Suporty Calibrações. Exibir a taxa de acerto na identificação de discrepâncias nas datas de calibração, vencimento e numeração dos laudos, demonstrando a eliminação de erros humanos em comparação ao processo manual anterior.)}

\section{Impacto Operacional no Acompanhamento do Farol}

\textcolor{red}{(Discutir aqui a validação do mecanismo em tempo real da projeção de instrumentos vincendos (Farol) e da guarda centralizada de evidências para instrumentos extraviados ou descartados, analisando a redução do esforço operacional dos instrumentistas e a melhoria na conformidade para auditorias de qualidade.)}
